% ============================================================================
\section{Scale Control Shifts the Variable-Depth Curve}
\label{sec:depth}

We next ask whether scale control changes how models use recurrent depth at inference time.
The result is not a generic speedup claim: a model can look fast simply because it is already $K$-invariant.
The stronger comparison is the fixed-depth PPL--compute frontier.

\begin{table}[t]
\centering
\small
\caption{\textbf{Scale-controlled per-loop models improve with depth and stay lower-PPL.}
Variable-depth inference on per-loop-loss models, mean PPL over 3 seeds. The $K=1$ and $K=4$ columns use the fixed-depth variable-depth slice harness; the dynamic-loop column summarizes a separately calibrated timed-slice halting protocol detailed in Appendix~\ref{app:halt-protocol}. Terminal-only variable-depth baselines (129M) are in Appendix~\ref{app:variable-depth}.}
\label{tab:depth}
\begin{tabular}{llrrrr}
\toprule
\textbf{Readout} & \textbf{Model} & $K=1$ & $K=4$ & $\Delta$PPL & Dynamic avg.\ loops \\
\midrule
RMSNorm & 44M & 5.58 & 5.59 & $+0.01$ & 1.00 \\
RMSNorm & 129M & 4.98 & 4.97 & $-0.01$ & 1.00 \\
Raw & 44M & 5.14 & 4.94 & $-0.20$ & 2.16 \\
Raw & 129M & 4.68 & 4.55 & $-0.13$ & 1.76 \\
Final-only norm & 44M & 5.20 & 5.00 & $-0.20$ & 1.78 \\
Final-only norm & 129M & 4.70 & 4.55 & $-0.15$ & 1.88 \\
Norm penalty & 44M & 5.22 & 5.00 & $-0.22$ & 2.60 \\
Norm penalty & 129M & 4.67 & 4.53 & $-0.14$ & 2.56 \\
\bottomrule
\end{tabular}
\end{table}

RMSNorm per-loop models are nearly $K$-invariant: at both sizes they can halt at $K=1$ under a 1\% PPL budget.
That produces high apparent speedup, but the model has little measurable use for additional loops.
This is the behavior Lemma~\ref{lemma:expansion} predicts at large scale: with final-loop norms in the tens of thousands, angular updates shrink like $1/s_k$, so additional loops barely change the normalized state the readout sees.
Norm drift therefore does not show up as a worse $K=4$ loss; its cost is the depth dimension itself.
Raw, final-only norm, and norm penalty all recover depth use: they remain usable at $K=1$ but improve with more loops.

\paragraph{High halting speedup can reflect $K$-invariance.}
We also benchmark a sequence-level logit-margin halting rule, calibrated to keep dynamic perplexity within $1\%$ of always running $K=4$.
The full dynamic-halting table is in Appendix~\ref{app:halt-protocol}; the result shows why a single speedup ratio is not the right endpoint for this paper.
The RMSNorm model exits at $K=1$ for every sequence and therefore appears fastest, but Table~\ref{tab:depth} shows that this happens because the model is already nearly $K$-invariant.
The scale-controlled rows spend more loops because additional loops actually improve perplexity.

The Pareto-relevant comparison is therefore the joint PPL--throughput frontier (Figure~\ref{fig:pareto}), not the speedup ratio in isolation.
For example, raw readout at $K=1$ has the same measured throughput as RMSNorm at $K=1$ but substantially lower PPL.

\begin{figure}[t]
\centering
\includegraphics[width=0.78\linewidth]{figures/pareto_compute_quality.pdf}
\caption{\textbf{Compute--quality frontier on per-loop-loss 129M models.}
Each curve sweeps $K\in\{1,2,3,4\}$.
Raw, final-only norm, and norm penalty occupy the lower-PPL region; RMSNorm is above them at every measured operating point.}
\label{fig:pareto}
\end{figure}

At equal throughput ($K=1$, about 80k tokens/s), raw readout is 0.30 PPL lower than RMSNorm.
At full depth ($K=4$, about 21k tokens/s), the gap is 0.42 PPL.
The conclusion is a compute--quality claim, not a higher-speedup claim: in this benchmark, scale-controlled per-loop variants preserve usable exits while shifting the measured frontier below the $K$-invariant RMSNorm baseline.

\paragraph{A simple practical intervention: scale-visibility penalty.}
For practitioners who want to keep normalized readouts, the norm penalty is the most direct intervention in our study.
The implementation is a one-line auxiliary loss,
$\lambda K^{-1}\sum_k \mathbb{E}\|H_k\|_{\rms}^2$, using hidden states already computed during the forward pass; its overhead is a reduction over activations, negligible compared with another recurrent block.
Table~\ref{tab:interventions-main} identifies it as an explicit scale-visible objective, Table~\ref{tab:main} shows that it collapses final-loop norms to the same range as raw readouts, and Table~\ref{tab:depth} shows that it preserves usable exits while recovering depth-dependent PPL improvement.
We do not claim this penalty is optimal.
Its role is simpler: it is a drop-in baseline that separates scale control from raw-readout confounds and is worth including whenever normalized readouts are kept but recurrent scale remains active.

Final-only normalization gives the complementary interface result: it uses raw intermediate readouts to control recurrence, but returns to a normalized final readout for the endpoint used by a fixed-depth model.
Together these controls make the conclusion less about one preferred readout and more about the underlying requirement.
Adaptive-depth looped models need trained exits and a way to control active recurrent scale.
Raw readout, an explicit penalty, and scale-removing recurrence are different engineering points on that same requirement; Appendix~\ref{app:mitigations} reports additional mitigation-style controls.
